\documentclass{article}
\usepackage[T1]{fontenc}
\usepackage{lmodern}
\usepackage{graphicx}
\usepackage{amsmath,amssymb}
\usepackage[a4paper,margin=2cm]{geometry}
\usepackage[absolute,overlay]{textpos}
\setlength{\TPHorizModule}{1mm}
\setlength{\TPVertModule}{1mm}
\usepackage[indonesian]{babel}
\usepackage{hyperref}
\setlength{\emergencystretch}{2em}
% unit: Halaman 41

\begin{document}

% === Halaman 41 (llm) ===

\section*{8. Sistem kriptografi (cryptosystem)}

Sistem kriptografi adalah $quintuple$ $(\mathcal{E}, \mathcal{D}, \mathcal{M}, \mathcal{K}, \mathcal{C})$
\begin{itemize}
    \item $\mathcal{M}$ adalah himpunan plainteks
    \item $\mathcal{K}$ adalah himpunan kunci
    \item $\mathcal{C}$ adalah himpunan cipherteks
    \item $\mathcal{E}$ adalah himpunan fungsi enkripsi $E: \mathcal{M} \times \mathcal{K} \to \mathcal{C}$
    \item $\mathcal{D}$ adalah himpunan fungsi dekripsi $D: \mathcal{C} \times \mathcal{K} \to \mathcal{M}$
\end{itemize}

\end{document}
